\chapter{Loadcases}
\label{chap:loadcases}

A loadcase defines which analysis is executed on the previously defined model. The model describes what exists: nodes, elements, surfaces, sets, materials, sections, features, loads, supports, constraints, and fields. The loadcase describes which equation is assembled, which collectors are active, how constraints are handled, which solver is used, and which analysis-specific parameters are applied.

\section{General syntax}

\begin{syntax}
*LOADCASE, TYPE=<type>, NAME=<optional-name>
  <loadcase commands>
\end{syntax}

\key{TYPE} is required. \key{NAME} is optional and is used for readability in the input deck and result output. Loadcases are executed in input order. Result blocks are written in the same order.

\section{Common equation structure}

Linear loadcases use global matrices assembled from element contributions. The most common operators are stiffness \(K\), mass \(M\), damping \(C\), and geometric stiffness \(K_G\). Active load collectors produce a right-hand side \(f\). Active supports and constraints produce equations
\[
C_c u = d .
\]
The selected constraint method then either reduces the problem with \(u=u_p+Tq\) or augments it with Lagrange multipliers. The loadcase formulas below are written in their physical full-space form first; the actual solved system may be the reduced or augmented form described in Chapter~\ref{chap:analysis-controls}.

\section{Linear Static}
\label{sec:loadcase-linear-static}

\val{LINEARSTATIC} solves small-displacement linear static equilibrium. It is the standard analysis for displacements, reactions, strains, stresses, and static load paths.

\subsection{Governing equation}

The physical equation is
\[
K u = f .
\]
After Nullspace constraint reduction, FEMaster solves
\[
T^T K T q = T^T(f-Ku_p),
\qquad
u=u_p+Tq .
\]
With Lagrange constraints, the static system is
\[
\begin{bmatrix}
K & C_c^T\\
C_c & 0
\end{bmatrix}
\begin{bmatrix}
u\\
\lambda
\end{bmatrix}
=
\begin{bmatrix}
f\\
d
\end{bmatrix}
\]
up to possible numerical regularization.

\subsection{Typical DSL}

\begin{exampledeck}
*LOADCASE, TYPE=LINEARSTATIC, NAME=static-load
  *SUPPORTS
  FIXED
  *LOADS
  FORCE
  *SOLVER, DEVICE=CPU, METHOD=DIRECT
  *CONSTRAINTMETHOD, TYPE=NULLSPACE
\end{exampledeck}

Common commands are \cmd{SUPPORTS}, \cmd{LOADS}, \cmd{SOLVER}, and \cmd{CONSTRAINTMETHOD}. Optional controls include \cmd{INERTIARELIEF}, \cmd{REBALANCELOADS}, \cmd{REQUESTSTIFFNESS}, and \cmd{CONSTRAINTSUMMARY}.

\section{Linear Static Topology}
\label{sec:loadcase-linear-static-topo}

\val{LINEARSTATICTOPO} solves a static problem whose element stiffness can be modified by topology fields. It is useful for topology studies, density-based parameter sweeps, and orientation-dependent stiffness experiments.

\subsection{Governing equation}

The equation is
\[
K(\rho,\theta)u=f .
\]
\(\rho\) is an element density-like field and \(\theta\) is an element orientation-like field. The density field is typically used through a penalized stiffness contribution
\[
K_e^\mathrm{eff} = \rho_e^p K_e ,
\]
where \(p\) is supplied by \cmd{TOPOEXPONENT}. Low densities reduce stiffness contribution; high densities approach the full element stiffness. The orientation field is used by element/material formulations that evaluate local orientation data.

\subsection{Typical DSL}

\begin{syntax}
*LOADCASE, TYPE=LINEARSTATICTOPO, NAME=<name>
  *TOPODENSITY
  <element-field>
  *TOPOORIENT
  <element-field>
  *TOPOEXPONENT
  <p>
\end{syntax}

\cmd{TOPODENSITY} requires an element field with one component. \cmd{TOPOORIENT} requires an element field with three components. The static controls from \val{LINEARSTATIC} are used in the same way.

The compliance, density sensitivity, and orientation sensitivity derivations are collected in Appendix~\ref{chap:appendix-derivations}.

\section{Eigenfrequency}
\label{sec:loadcase-eigenfrequency}

\val{EIGENFREQ} computes natural frequencies and mode shapes of the linearized model. Loads are not used because the free-vibration problem is defined by stiffness, mass, and boundary conditions.

\subsection{Governing equation}

The full-space eigenproblem is
\[
K\phi_i = \omega_i^2 M\phi_i .
\]
\(\phi_i\) is the mode shape and \(\omega_i\) is the circular frequency. The frequency in Hertz is
\[
f_i = \frac{\omega_i}{2\pi}.
\]
With Nullspace constraints, the reduced problem is
\[
T^TKT\psi_i = \omega_i^2 T^TMT\psi_i,
\qquad
\phi_i=T\psi_i .
\]

\subsection{Typical DSL}

\begin{syntax}
*LOADCASE, TYPE=EIGENFREQ, NAME=<name>
  *SUPPORTS
  <support-collector>
  *NUMEIGENVALUES
  <count>
\end{syntax}

A free-free model can contain rigid-body modes with zero or near-zero frequency. A supported model computes modes relative to the active supports and constraints.

\val{EIGENFREQ} uses Nullspace constraints internally. \cmd{CONSTRAINTMETHOD} is not a valid command for this loadcase.

\section{Linear Buckling}
\label{sec:loadcase-linear-buckling}

\val{LINEARBUCKLING} computes ideal linear buckling factors and buckling modes around a preloaded reference state. It first solves a static preload problem, then assembles geometric stiffness from the resulting stresses, and finally solves a generalized eigenvalue problem.

\subsection{Governing equations}

The preload solve is
\[
K u_0 = f_\mathrm{pre}.
\]
With constraints, FEMaster solves the reduced preload equation
\[
A q_\mathrm{pre}=b,\qquad A=T^TKT,\qquad b=T^T(f_\mathrm{pre}-Ku_p).
\]
From the preload stresses, FEMaster builds the geometric stiffness \(K_G\) and reduces it to
\[
B=T^T K_G T .
\]
The buckling eigenproblem is solved as
\[
A\phi_i=\lambda_i(-B)\phi_i .
\]
\(\lambda_i\) is the buckling factor. For the preload pattern \(f_\mathrm{pre}\), the ideal linear critical load for mode \(i\) is
\[
f_{\mathrm{crit},i} = \lambda_i f_\mathrm{pre}.
\]

\subsection{Sigma}

\cmd{SIGMA} controls the buckling eigenvalue shift. It should be set carefully because it guides the part of the spectrum returned by the eigensolver. If \cmd{SIGMA} is zero, FEMaster estimates a shift automatically from the preload displacement with the Rayleigh quotient
\[
\lambda_\mathrm{raw}
=
\frac{q_\mathrm{pre}^T A q_\mathrm{pre}}
       {q_\mathrm{pre}^T(-B)q_\mathrm{pre}} .
\]
When this estimate is valid and positive, the generated shift is \(\lambda_\mathrm{raw}/10^6\). Otherwise a diagonal-ratio fallback is used; if that also fails, \(10^{-12}\) is used.

\subsection{Typical DSL}

\begin{syntax}
*LOADCASE, TYPE=LINEARBUCKLING, NAME=<name>
  *SUPPORTS
  <support-collector>
  *LOADS
  <preload-collector>
  *NUMEIGENVALUES
  <count>
  *SIGMA
  <shift>
\end{syntax}

\cmd{LOADS} defines the preload. \cmd{REQUESTSTIFFNESS} and \cmd{REQUESTSTGEOM} can be used to write elastic and geometric stiffness matrices for diagnostics.

\val{LINEARBUCKLING} uses Nullspace constraints internally. \cmd{CONSTRAINTMETHOD} is not a valid command for this loadcase.

\section{Nonlinear Static}
\label{sec:loadcase-nonlinear-static}

\val{NONLINEARSTATIC} solves quasi-static equilibrium by applying the active load in increments and iterating with a tangent stiffness matrix in each increment. It is the only loadcase that assembles \cmd{CONTACT}. It also uses the nonlinear element tangent/internal-force path where available.

\subsection{Governing equation}

At a load factor \(\lambda\), the physical equilibrium equation is
\[
f_\mathrm{int}(u) = \lambda f_\mathrm{total}.
\]
FEMaster stores a reference position field and updates current nodal positions from the total displacement. Translational coordinates are updated as
\[
x = X + u ,
\]
while rotational entries are treated as total generalized displacement components. Linear constraints are enforced with the Nullspace representation
\[
u = u_p + Tq .
\]
The reduced Newton residual is
\[
\hat r(q)=T^T(\lambda f_\mathrm{total}-f_\mathrm{int}(u_p+Tq)).
\]
The reduced tangent equation solved in each Newton iteration is
\[
T^T K_t T\,\Delta q = \hat r .
\]
The current implementation applies the full correction \(u \leftarrow u + T\Delta q\).

\subsection{Convergence and stepping}

The relative residual reported in the nonlinear iteration log is the reduced residual RMS normalized by the larger of the reduced external-force RMS, the reduced internal-force RMS, and 1. The increment is accepted when this value is less than or equal to \key{TOL}. If convergence fails within \key{MAXITER}, FEMaster restores the last accepted state and cuts the load increment back. If convergence is fast, the next increment is allowed to grow.

\cmd{NONLINEAR} controls the initial increment count, maximum iterations, residual tolerance, and optional zero-row tangent regularization. Detailed syntax is given in Chapter~\ref{chap:analysis-controls}.

\subsection{Typical DSL}

\begin{syntax}
*LOADCASE, TYPE=NONLINEARSTATIC, NAME=<name>
  *SUPPORTS
  <support-collector>
  *LOADS
  <load-collector>
  *SOLVER, DEVICE=CPU, METHOD=DIRECT
  *CONSTRAINTMETHOD, TYPE=NULLSPACE
  *NONLINEAR, INCREMENTS=20, MAXITER=25, TOL=1e-8
\end{syntax}

\val{NONLINEARSTATIC} currently requires \val{DIRECT} solver method and \val{NULLSPACE} constraints. \cmd{REQUESTSTIFFNESS} writes the final tangent stiffness \(K_t\) and final reduced tangent \(T^TK_tT\). Contact definitions are global model definitions, not loadcase commands; they contribute only when this loadcase is executed.

\section{Linear Transient}
\label{sec:loadcase-linear-transient}

\val{LINEARTRANSIENT} solves the linear dynamic response over time with implicit Newmark integration. The structure remains linear; contact, plasticity, and large deformation effects are not included in this loadcase.

\subsection{Governing equation}

The physical equation is
\[
M\ddot u(t)+C\dot u(t)+Ku(t)=f(t).
\]
With Nullspace constraints, this becomes
\[
M_r\ddot q(t)+C_r\dot q(t)+K_rq(t)=r(t),
\]
with
\[
M_r=T^TMT,\qquad C_r=T^TCT,\qquad K_r=T^TKT .
\]
Loads with amplitudes are assembled as
\[
f(t)=\sum_j a_j(t)f_j .
\]
For Rayleigh damping,
\[
C=\alpha M+\beta K .
\]

\subsection{Typical DSL}

\begin{syntax}
*LOADCASE, TYPE=LINEARTRANSIENT, NAME=<name>
  *SUPPORTS
  <support-collector>
  *LOADS
  <load-collector>
  *TIME
  <t_start>, <t_end>, <dt>
  *NEWMARK
  <beta>, <gamma>
  *DAMPING, TYPE=RAYLEIGH
  <alpha>, <beta>
  *WRITEEVERY, TYPE=STEPS
  <n>
\end{syntax}

The transient-specific keywords \cmd{TIME}, \cmd{NEWMARK}, \cmd{DAMPING}, \cmd{WRITEEVERY}, and \cmd{INITIALVELOCITY} are described in Chapter~\ref{chap:analysis-controls}.

\val{LINEARTRANSIENT} uses Nullspace constraints internally. \cmd{CONSTRAINTMETHOD} is not a valid command for this loadcase.

\section{Typical command overview}

\begin{longtable}{@{}p{0.24\linewidth}p{0.66\linewidth}@{}}
\toprule
\textbf{Loadcase} & \textbf{Typical commands} \\
\midrule
\endhead
\val{LINEARSTATIC} & \cmd{SUPPORTS}, \cmd{LOADS}, \cmd{SOLVER}, \cmd{CONSTRAINTMETHOD}, optional \cmd{INERTIARELIEF}, \cmd{REBALANCELOADS}, \cmd{REQUESTSTIFFNESS}, \cmd{CONSTRAINTSUMMARY}. \\
\val{LINEARSTATICTOPO} & Linear-static commands plus \cmd{TOPODENSITY}, \cmd{TOPOORIENT}, \cmd{TOPOEXPONENT}. \\
\val{NONLINEARSTATIC} & \cmd{SUPPORTS}, \cmd{LOADS}, \cmd{SOLVER}, \cmd{CONSTRAINTMETHOD} with \val{NULLSPACE}, \cmd{NONLINEAR}, optional \cmd{REQUESTSTIFFNESS}, \cmd{CONSTRAINTSUMMARY}. \\
\val{EIGENFREQ} & \cmd{SUPPORTS}, \cmd{NUMEIGENVALUES}. \\
\val{LINEARBUCKLING} & \cmd{SUPPORTS}, \cmd{LOADS}, \cmd{NUMEIGENVALUES}, optional \cmd{SIGMA}, \cmd{REQUESTSTIFFNESS}, \cmd{REQUESTSTGEOM}. \\
\val{LINEARTRANSIENT} & \cmd{SUPPORTS}, \cmd{LOADS}, \cmd{SOLVER}, \cmd{TIME}, \cmd{NEWMARK}, \cmd{DAMPING}, \cmd{WRITEEVERY}, \cmd{INITIALVELOCITY}. \\
\bottomrule
\end{longtable}
